% =====================================================================
%  مقدمة الناظم (٢): الأبيات ٢٠–٤٣ — القراء السبعة ورواتهم
% =====================================================================
\chapter{مقدمة الناظم (٢): القراء السبعة ورواتهم}

\begin{tanbih}[فكرة هذا القسم]
يعرّفك الناظم هنا بالقراء السبعة، ولكل قارئ راويان. حفظ هذه الأسماء بترتيبها \textbf{مهم جدًّا}؛ لأن رموز «أبجد» مبنية على هذا الترتيب نفسه.
ولوّنا اسم كل قارئ وراوييه بلونه في الأبيات.
\end{tanbih}

\bayt{20}{جَزَى اللهُ بِالْخَيْرَاتِ عَنَّا أَئِمَّةً}{لَنَا نَقَلُوا الْقُرْآنَ عَذْبًا وَسَلْسَلَا}
\begin{mufradat}
\kalima{عذبًا وسلسلًا}{حلوًا سهلًا على اللسان.}
\end{mufradat}
\maana دعاء للأئمة الذين نقلوا إلينا القرآن نقلًا صحيحًا متقنًا، سهلًا عذبًا كما أُنزل.

\bayt{21}{فَمِنْهُمْ بُدُورٌ سَبْعَةٌ قَدْ تَوَسَّطَتْ}{سَمَاءَ الْعُلَى وَالْعَدْلِ زُهْرًا وَكُمَّلَا}
\begin{mufradat}
\kalima{بدور}{جمع بدر، والمراد القراء السبعة.}
\kalima{زُهرًا}{مضيئة مشرقة.}
\kalima{كُمَّلًا}{جمع كامل.}
\end{mufradat}
\maana من هؤلاء الأئمة سبعة كالبدور، توسطوا سماء الرفعة والعدالة، مشرقين كاملين.

\bayt{22}{لَهَا شُهُبٌ عَنْهَا اسْتَنَارَتْ فَنَوَّرَتْ}{سَوَادَ الدُّجَى حَتَّى تَفَرَّقَ وَانْجَلَى}
\begin{mufradat}
\kalima{شُهُب}{جمع شهاب، وهو النجم المضيء، والمراد الرواة.}
\kalima{الدُّجى}{الظلام.}
\end{mufradat}
\maana ولهؤلاء البدور (القراء) نجوم (رواة) أخذوا النور منهم، فأضاؤوا ظلام الجهل حتى انكشف.

\bayt{23}{وَسَوْفَ تَرَاهُمْ وَاحِدًا بَعْدَ وَاحِدٍ}{مَعَ اثْنَيْنِ مِنْ أَصْحَابِهِ مُتَمَثِّلَا}
\begin{mufradat}
\kalima{متمثلًا}{ظاهرًا ماثلًا أمامك.}
\end{mufradat}
\maana سأذكر لك القراء واحدًا واحدًا، ومع كل قارئ راويين من أصحابه.

\bayt{24}{تَخَيَّرَهُمْ نُقَّادُهُمْ كُلَّ بَارِعٍ}{وَلَيْسَ عَلَى قُرْآنِهِ مُتَأَكِّلَا}
\begin{mufradat}
\kalima{نُقّادهم}{العلماء الذين ميّزوا الرواة وانتقوهم.}
\kalima{بارع}{فائق في علمه.}
\kalima{متأكّلًا}{يجعل القرآن وسيلة للتكسّب.}
\end{mufradat}
\maana اختار العلماءُ النقادُ هؤلاء الرواة؛ فكل واحد منهم بارع متقن، ولا يتخذ القرآن وسيلة للأكل والكسب.

\section{نافع المدني وراوياه}
\bayt{25}{فَأَمَّا الْكَرِيمُ السِّرِّ فِي الطِّيبِ \qname{Nafi}{نَافِعٌ}}{فَذَاكَ الَّذِي اخْتَارَ الْمَدِينَةَ مَنْزِلَا}
\begin{mufradat}
\kalima{الكريم السر}{طاهر الباطن.}
\kalima{في الطِّيب}{يُروى أنه كان إذا تكلم شُمّت من فمه رائحة المسك، فسئل عن ذلك فقال: رأيت النبي ﷺ في المنام يقرأ في فيّ.}
\end{mufradat}
\maana القارئ الأول: \Nafi{} بن عبد الرحمن بن أبي نُعيم، اختار المدينة المنورة دارًا ومقامًا (ت ١٦٩ هـ).

\bayt{26}{\qname{Nafi}{وَقَالُونُ} عِيسَى ثُمَّ عُثْمَانُ \qname{Nafi}{وَرْشُهُمْ}}{بِصُحْبَتِهِ الْمَجْدَ الرَّفِيعَ تَأَثَّلَا}
\begin{mufradat}
\kalima{تأثّلا}{اكتسبا وأصّلا.}
\end{mufradat}
\maana راويا \Nafi: \Qalun{} واسمه \textbf{عيسى} (ت ٢٢٠ هـ)، ثم \Warsh{} واسمه \textbf{عثمان} (ت ١٩٧ هـ)، وقد نالا المجد الرفيع بصحبتهما \Nafi.

\section{ابن كثير المكي وراوياه}
\bayt{27}{وَمَكَّةُ عَبْدُ اللهِ فِيهَا مُقَامُهُ}{هُوَ \qname{IbnKathir}{ابْنُ كَثِيرٍ} كَاثِرُ الْقَوْمِ مُعْتَلَى}
\begin{mufradat}
\kalima{كاثر القوم}{غالبهم وفائقهم (وفيه تلاعب لطيف باسم «كثير»).}
\kalima{مُعتلى}{عاليًا رفيعًا.}
\end{mufradat}
\maana القارئ الثاني: \textbf{عبد الله} \IbnKathir{} المكي، أقام بمكة، وفاق قومه علمًا ورفعة (ت ١٢٠ هـ).

\bayt{28}{رَوَى أَحْمَدُ \qname{IbnKathir}{الْبَزِّي} لَهُ وَمُحَمَّدٌ}{عَلَى سَنَدٍ وَهْوَ الْمُلَقَّبُ \qname{IbnKathir}{قُنْبُلَا}}
\begin{mufradat}
\kalima{على سند}{أي رويا عنه بواسطة رجال، لا مباشرة.}
\end{mufradat}
\maana راويا \IbnKathir: \Bazzi{} واسمه \textbf{أحمد} (ت ٢٥٠ هـ)، و\Qunbul{} واسمه \textbf{محمد} (ت ٢٩١ هـ)، وكلاهما روى عنه بواسطة.

\section{أبو عمرو البصري وراوياه}
\bayt{29}{وَأَمَّا الْإِمَامُ الْمَازِنِيُّ صَرِيحُهُمْ}{\qname{AbuAmr}{أَبُو عَمْرٍو الْبَصْرِي} فَوَالِدُهُ الْعَلَا}
\begin{mufradat}
\kalima{المازني}{نسبة إلى بني مازن من تميم.}
\kalima{صريحهم}{خالص النسب العربي فيهم.}
\end{mufradat}
\maana القارئ الثالث: \AbuAmr{} البصري، عربي صريح النسب من بني مازن، وأبوه اسمه \textbf{العلاء} (ت ١٥٤ هـ). ولذلك سيسميه الناظم كثيرًا «ابن العلاء» و«فتى العلا».

\bayt{30}{أَفَاضَ عَلَى يَحْيَى الْيَزِيدِيِّ سَيْبَهُ}{فَأَصْبَحَ بِالْعَذْبِ الْفُرَاتِ مُعَلَّلَا}
\begin{mufradat}
\kalima{سَيْبَه}{عطاءه، أي علمه.}
\kalima{العذب الفرات}{الماء الحلو الشديد العذوبة.}
\kalima{معلَّلًا}{مسقيًّا مرة بعد مرة.}
\end{mufradat}
\maana أفاض \AbuAmr{} علمه على تلميذه \textbf{يحيى اليزيدي}، فارتوى منه، وعن اليزيدي أخذ الراويان المذكوران بعده.

\bayt{31}{أَبُو عُمَرَ \qname{AbuAmr}{الدُّورِي} وَصَالِحُهُمْ أَبُو}{شُعَيْبٍ هُوَ \qname{AbuAmr}{السُّوسِيُّ} عَنْهُ تَقَبَّلَا}
\maana راويا \AbuAmr{} (بواسطة اليزيدي): \qname{AbuAmr}{الدوري} حفص بن عمر، كنيته \textbf{أبو عُمَر} (ت ٢٤٦ هـ)، و\Susi{} صالح بن زياد، كنيته \textbf{أبو شعيب} (ت ٢٦١ هـ).

\section{ابن عامر الشامي وراوياه}
\bayt{32}{وَأَمَّا دِمَشْقُ الشَّامِ دَارُ \qname{IbnAmir}{ابْنِ عَامِرٍ}}{فَتِلْكَ بِعَبْدِ اللهِ طَابَتْ مُحَلَّلَا}
\begin{mufradat}
\kalima{مُحَلَّلًا}{منزلًا ومكانَ حلول.}
\end{mufradat}
\maana القارئ الرابع: \textbf{عبد الله} \IbnAmir{} اليَحْصُبي، داره دمشق، وطابت به منزلًا (ت ١١٨ هـ). وهو أعلى القراء السبعة سندًا.

\bayt{33}{\qname{IbnAmir}{هِشَامٌ} وَعَبْدُ اللهِ وَهْوَ انْتِسَابُهُ}{\qname{IbnAmir}{لِذَكْوَانَ} بِالْإِسْنَادِ عَنْهُ تَنَقَّلَا}
\maana راويا \IbnAmir: \Hisham{} بن عمار (ت ٢٤٥ هـ)، و\textbf{عبد الله} بن أحمد \IbnDhakwan{} — نُسب إلى جده ذكوان (ت ٢٤٢ هـ). وكلاهما روى عنه بواسطة.

\section{القراء الكوفيون الثلاثة}
\bayt{34}{وَبِالْكُوفَةِ الْغَرَّاءِ مِنْهُمْ ثَلَاثَةٌ}{أَذَاعُوا فَقَدْ ضَاعَتْ شَذًا وَقَرَنْفُلَا}
\begin{mufradat}
\kalima{الغرّاء}{البيضاء المشهورة.}
\kalima{أذاعوا}{نشروا القراءة.}
\kalima{ضاعت}{فاحت رائحتها.}
\kalima{شذًا}{الرائحة الطيبة.}
\end{mufradat}
\maana وفي الكوفة ثلاثة من السبعة نشروا القرآن، ففاحت الكوفة بعلمهم كرائحة الطيب والقرنفل. وهؤلاء الثلاثة هم «الكوفيون» ورمزهم \rzlg{ث}.

\bayt{35}{فَأَمَّا أَبُو بَكْرٍ وَ\qname{Asim}{عَاصِمٌ} اسْمُهُ}{\qname{Asim}{فَشُعْبَةُ} رَاوِيهِ الْمُبَرِّزُ أَفْضَلَا}
\begin{mufradat}
\kalima{المبرِّز}{المتقدم الفائق.}
\end{mufradat}
\maana القارئ الخامس: \Asim{} بن أبي النَّجود، كنيته \textbf{أبو بكر} (ت ١٢٧ هـ). وراويه الأول \Shuba، المتقدم الفاضل.

\bayt{36}{وَذَاكَ ابْنُ عَيَّاشٍ أَبُو بَكْرٍ الرِّضَا}{وَ\qname{Asim}{حَفْصٌ} وَبِالْإِتْقَانِ كَانَ مُفَضَّلَا}
\maana \Shuba{} هو ابن عيّاش، وكنيته أيضًا \textbf{أبو بكر} (ت ١٩٣ هـ). والراوي الثاني \Hafs{} بن سليمان، وكان مفضّلًا بالإتقان (ت ١٨٠ هـ). وروايته هي المنتشرة اليوم في أكثر العالم الإسلامي.

\bayt{37}{وَ\qname{Hamza}{حَمْزَةُ} مَا أَزْكَاهُ مِنْ مُتَوَرِّعٍ}{إِمَامًا صَبُورًا لِلْقُرَانِ مُرَتِّلَا}
\maana القارئ السادس: \Hamza{} بن حبيب الزيّات، ما أطهره من ورِع! كان إمامًا صبورًا مرتّلًا للقرآن (ت ١٥٦ هـ).

\bayt{38}{رَوَى \qname{Hamza}{خَلَفٌ} عَنْهُ وَ\qname{Hamza}{خَلَّادٌ} الَّذِي}{رَوَاهُ سُلَيْمٌ مُتْقَنًا وَمُحَصَّلَا}
\maana راويا \Hamza: \Khalaf{} بن هشام (ت ٢٢٩ هـ) و\Khallad{} بن خالد (ت ٢٢٠ هـ)، وكلاهما أخذ عن \textbf{سُلَيم} بن عيسى عن \Hamza.

\bayt{39}{وَأَمَّا عَلِيٌّ \qname{Kisai}{فَالْكِسَائِيُّ} نَعْتُهُ}{لِمَا كَانَ فِي الْإِحْرَامِ فِيهِ تَسَرْبَلَا}
\begin{mufradat}
\kalima{تسربل}{لبس.}
\end{mufradat}
\maana القارئ السابع: \textbf{علي} بن حمزة، ولقبه \Kisai؛ لأنه أحرم في كساء (ت ١٨٩ هـ).

\bayt{40}{رَوَى لَيْثُهُمْ عَنْهُ \qname{Kisai}{أَبُو الْحَارِثِ} الرِّضَا}{وَحَفْصٌ هُوَ \qname{Kisai}{الدُّورِي} وَفِي الذِّكْرِ قَدْ خَلَا}
\begin{mufradat}
\kalima{ليثهم}{اسمه الليث بن خالد.}
\kalima{في الذكر قد خلا}{سبق ذكره (مع أبي عمرو).}
\end{mufradat}
\maana راويا \Kisai: \AbulHarith{} الليث بن خالد (ت ٢٤٠ هـ)، و\textbf{حفص الدوري}، وهو نفسه راوي \AbiAmr{} الذي سبق ذكره.

\begin{tanbih}[انتبه: الدوري راوٍ لقارئين]
الدوري (حفص بن عمر) روى عن \AbiAmr{} وروى أيضًا عن \Kisai. فإذا قيل «الدوري» في باب فانظر: هل الكلام عن \AbiAmr{} أم \Kisai؟ ولهذا جعلنا لكل منهما لون قارئه.
\end{tanbih}

\section{أنساب القراء وطرق الرواية}
\bayt{41}{\qname{AbuAmr}{أَبُو عَمْرِهِمْ} وَالْيَحْصَبِيُّ \qname{IbnAmir}{ابْنُ عَامِرٍ}}{صَرِيحٌ وَبَاقِيهِمْ أَحَاطَ بِهِ الْوَلَا}
\begin{mufradat}
\kalima{صريح}{عربي خالص النسب.}
\kalima{الوَلا}{الولاء، أي أنهم من الموالي (أصولهم غير عربية).}
\end{mufradat}
\maana \AbuAmr{} و\IbnAmir{} عربيان صريحا النسب، والباقون من الموالي. (وهذا لا ينقص من قدرهم شيئًا، فالفضل بالعلم والتقوى.)

\bayt{42}{لَهُمْ طُرُقٌ يُهْدَى بِهَا كُلُّ طَارِقٍ}{وَلَا طَارِقٌ يُخْشَى بِهَا مُتَمَحِّلَا}
\begin{mufradat}
\kalima{طُرُق}{طرق الرواية والأسانيد.}
\kalima{طارق}{قاصد (من يأتي ليلًا).}
\kalima{ولا طارقٌ يُخشى}{ولا داهية أو ضلال يُخاف منه.}
\kalima{متمحِّلًا}{ماكرًا كائدًا.}
\end{mufradat}
\maana لهؤلاء الرواة طرق (أسانيد) يهتدي بها كل طالب، ولا يُخشى فيها ضلال.

\begin{qaida}[ما الفرق بين القراءة والرواية والطريق؟]
\textbf{القراءة}: ما يُنسب إلى الإمام (مثل: قراءة \Nafi).\\
\textbf{الرواية}: ما يُنسب إلى الراوي عن الإمام (مثل: رواية \Qalun{} عن \Nafi).\\
\textbf{الطريق}: ما يُنسب إلى من أخذ عن الراوي وإن نزل (مثل: طريق أبي نشيط عن \Qalun، وهو طريق الشاطبية).\\
والشاطبية نَظْمٌ لكتاب «التيسير» للداني، فكل راوٍ فيها من طريق واحد غالبًا.
\end{qaida}

\bayt{43}{وَهُنَّ اللَّوَاتِي لِلْمُوَاتِي نَصَبْتُهَا}{مَنَاصِبَ فَانْصَبْ فِي نِصَابِكَ مُفْضِلَا}
\begin{mufradat}
\kalima{المُواتي}{الموافق المطاوع، أي الطالب المجتهد.}
\kalima{مناصب}{علامات منصوبة يُهتدى بها.}
\kalima{فانْصَبْ}{اجتهد واتعب.}
\kalima{نِصابك}{أصلك وحظك.}
\kalima{مُفضِلًا}{محسنًا.}
\end{mufradat}
\maana هذه الطرق نصبتُها علامات للطالب الموافق، فاجتهد في طلب هذا العلم محسنًا.
